% Drop-in body only. Merge preamble macros and bibliography entries first.
\section{Wireless-Aware Joint Bandwidth and Posterior-Retention Allocation}
\label{sec:fedkls-wireless}
\subsection{Wireless System and Synchronized Timing}
\label{fkl:sec:system}
We consider one server and $K$ participating clients indexed by $\mathcal K=\{1,\ldots,K\}$. The common neural network has $d\ge2$ scalar parameters. Client $k$ holds $n_k>0$ samples and has aggregation weight $\pi_k=n_k/\sum_j n_j$. As in FedKLS~\cite{fedkls}, each client trains the full local model and subsequently transmits selected posterior mean--variance pairs. Sparsity therefore changes the message, not the local training architecture. The proposed extension allocates bandwidth and retention based on current wireless conditions and posterior changes; it does not replace the underlying Bayesian learning algorithm.

\subsubsection{Orthogonal Digital Links and Channel Model}
Clients upload concurrently over nonoverlapping frequency resources. The allocated bandwidths satisfy
\begin{equation}
 B_{k,t}>0,\qquad \sum_{k\in\mathcal K}B_{k,t}\le B,
 \label{fkl:eq:bandwidth-budget}
\end{equation}
where $B$ is the total uplink bandwidth. Orthogonality suppresses intra-cell interference in the analytical model. This frequency-allocation abstraction is compatible with, but does not require, a particular multicarrier waveform. In contrast, the source scheduling framework explicitly chooses OFDMA~\cite[Sec.~II-C]{scheduling}.

A proposed block-fading model separates path gain and small-scale fading:
\begin{equation}
 H_{k,t}=\sqrt{\beta_k}\,h_{k,t},\qquad
 \beta_k=\beta_0\left(\frac{\ell_k}{\ell_0}\right)^{-\alpha},
 \label{fkl:eq:channel}
\end{equation}
where $\ell_k\ge\ell_0>0$ is the client--server distance, $\beta_0$ is the gain at reference distance $\ell_0$, and $\alpha$ is the path-loss exponent. For the initial simulation, choose $h_{k,t}\sim\mathcal{CN}(0,1)$ independently between rounds and keep the gain fixed over the current allocation and upload phase. This Rayleigh assumption is a modeling choice, not implied by the rate formula. Fixed client positions isolate persistent distance effects from round-to-round fading. Perfect current channel knowledge is assumed for the initial allocator.

The uplink uses the idealized rate model
\begin{equation}
 R_{k,t}^{\UL}=B_{k,t}c_{k,t},\qquad
 c_{k,t}=\log_2(1+\gamma_{k,t}^{\UL})>0.
 \label{fkl:eq:rate}
\end{equation}
Here $c_{k,t}$ is spectral efficiency, held fixed during the bandwidth decision. One physically consistent convention is a fixed transmit power spectral density $p_{k,t}^{\mathrm{PSD}}$ (W/Hz):
\begin{equation}
 \gamma_{k,t}^{\UL}
 =\frac{p_{k,t}^{\mathrm{PSD}}\beta_k|h_{k,t}|^2}{N_0},
 \qquad P_{k,t}=p_{k,t}^{\mathrm{PSD}}B_{k,t},
 \label{fkl:eq:snr-psd}
\end{equation}
where $N_0$ is noise power spectral density. The allocation problem does not optimize power. If a total-power limit $P_k^{\max}$ is to hold automatically without adding a bandwidth cap, a sufficient configuration is $p_{k,t}^{\mathrm{PSD}}B\le P_k^{\max}$. The actual power and energy must then be evaluated using the allocated $B_{k,t}$.

Under a different convention with fixed total transmit power, the rate would instead be
\begin{equation}
 R_{k,t}^{\UL}(B_{k,t})
 =B_{k,t}\log_2\!\left(1+
 \frac{P_{k,t}\beta_k|h_{k,t}|^2}{N_0B_{k,t}}\right).
 \label{fkl:eq:fixed-power-caution}
\end{equation}
Equation~\eqref{fkl:eq:fixed-power-caution} is not used in the derivations below. Holding both total power and SNR fixed while varying bandwidth would be inconsistent with this noise model.

\subsubsection{Posterior Payload and Transmission Time}
Let $b$ denote the bits used to represent one posterior mean or variance. Client $k$ transmits $m_{k,t}$ pairs, where
\begin{equation}
 m_{k,t}\in\{1,\ldots,d\},\qquad
 r_{k,t}=\frac{m_{k,t}}{d}\in(0,1].
 \label{fkl:eq:retention}
\end{equation}
Thus, every integer retained count is allowed; no preset percentage menu is introduced. The percentage keep ratio is $100r_{k,t}$, with unavoidable resolution $100/d$ percentage points.

With the submitted FedKLS bitmap representation~\cite[Sec.~II-B]{fedkls}, the uplink payload is
\begin{equation}
 L_{k,t}^{\UL}(m)=
 \begin{cases}
 d+2bm,&1\le m<d,\\
 2bd,&m=d.
 \end{cases}
 \label{fkl:eq:uplink-payload}
\end{equation}
The $d$-bit bitmap identifies sparse coordinates. At full retention, coordinate order is known and the bitmap is omitted. Additional headers, byte padding, and packet fragmentation are neglected, but the bitmap is retained. Payload $L$ is measured in bits, not in time. Accordingly,
\begin{equation}
 T_{k,t}^{\UL}(m_{k,t},B_{k,t})
 =\frac{L_{k,t}^{\UL}(m_{k,t})}{B_{k,t}c_{k,t}}.
 \label{fkl:eq:upload-time}
\end{equation}

The server sends the full global posterior by reliable digital multicast. Its payload is $2bd$ bits, counted once. Using a rate supported by the worst intended receiver, as in~\cite[Eqs.~(4)--(5)]{scheduling},
\begin{equation}
 T_t^{\DL}=\frac{2bd}{B^{\DL}\log_2\!\left(
 1+\min_{k\in\mathcal K}\gamma_{k,t}^{\DL}\right)}.
 \label{fkl:eq:downlink-time}
\end{equation}
Reducing uplink retention does not shorten this dense downlink. Reliable decoding means that successfully delivered values equal their transmitted finite-precision representation; no additive channel noise is injected directly into those decoded values in the initial model. Additional quantization, packet loss, and retransmissions are outside the initial formulation.

\subsubsection{Local Preparation and the Reporting Barrier}
The general local preparation duration is
\begin{align}
 T_{k,t}^{\local}
 &=T_{k,t}^{\train}+T_{k,t}^{\KL},\label{fkl:eq:local-time}\\
 T_{k,t}^{\train}
 &\simeq\frac{E_{k,t}n_kC_{\mathrm{FOLA}}}
 {f_{k,t}^{\mathrm{eff}}},\label{fkl:eq:train-time}
\end{align}
where $E_{k,t}$ is the number of local epochs, $C_{\mathrm{FOLA}}$ is the effective computation per processed example including the adopted posterior-curvature estimator, and $f_{k,t}^{\mathrm{eff}}$ is throughput in operations/s. Equation~\eqref{fkl:eq:train-time} adapts the operation-count latency model in~\cite[Eq.~(16)]{scheduling}; profiled durations may be substituted. The initial approximation $T_{k,t}^{\local}\simeq T_{k,t}^{\train}$ neglects KL computation and ranking time. Alternatively, computing a complete score ranking once permits its measured cost to be included without making local preparation depend on the allocated count.

The protocol is multicast, parallel local preparation, collection of current importance information and channel information, allocation, concurrent posterior uploads, and aggregation. All clients wait until the current information required by the allocator exists. A client finishing early therefore has pre-upload waiting time
\begin{equation}
 W_{k,t}^{\mathrm{pre}}
 =\max_iT_{i,t}^{\local}-T_{k,t}^{\local}.
 \label{fkl:eq:waiting}
\end{equation}
The per-client elapsed completion time is
\begin{equation}
 \begin{split}
 T_{k,t}^{\round}(m_{k,t},B_{k,t})
 ={}&T_t^{\DL}+\max_iT_{i,t}^{\local}\\
 &+T_{k,t}^{\UL}(m_{k,t},B_{k,t}).
 \end{split}
 \label{fkl:eq:round-time}
\end{equation}
This preserves the staged structure of~\cite[Eq.~(18)]{scheduling}. In that source, the client index specifies the sole selected uploader; here it specifies an individual completion time in a multi-client round. It does not denote independent federated clocks.

Channel-estimation, control-report, grant, packing, server-allocation, and aggregation durations are omitted initially. Their omission is a latency approximation, not permission to allocate using information before it is available. In particular, the proposed iterative price coordination occurs after the reporting barrier and before the main uploads.

\subsubsection{Prescribed Common Completion Time}
Let $T_t^{\round}$, without client index, be a prescribed common completion target. It and the current downlink and local durations are known when allocation begins. Define
\begin{equation}
 U_t=T_t^{\round}-T_t^{\DL}-\max_iT_{i,t}^{\local}>0.
 \label{fkl:eq:uplink-budget}
\end{equation}
The equality-based allocation enforces $T_{k,t}^{\round}=T_t^{\round}$ for all clients, or equivalently
\begin{equation}
 L_{k,t}^{\UL}(m_{k,t})=U_tB_{k,t}c_{k,t}.
 \label{fkl:eq:time-equality}
\end{equation}
Counts are chosen as integers, and bandwidth is then recovered as
\begin{equation}
 B_{k,t}=\frac{L_{k,t}^{\UL}(m_{k,t})}{U_tc_{k,t}}.
 \label{fkl:eq:bandwidth-recovery}
\end{equation}
This yields exactly $U_t$ seconds of idealized uploading and may leave some total bandwidth unused. It is not the policy that uses all bandwidth to minimize completion time. At a fixed bandwidth, flooring $(U_tB_{k,t}c_{k,t}-d)/(2b)$ only enforces a sparse packet-fitting inequality, not an exact time equality.


\subsection{Joint Bandwidth and Posterior-Retention Allocation}
\label{fkl:sec:problem}
\subsubsection{Exact Local Reconstruction Cost}
Let $q_{t,0}$ be the broadcast diagonal Gaussian posterior and $q_{k,t}$ the full local diagonal Gaussian approximation after training. All variances are strictly positive. The original FedKLS omission score~\cite[Sec.~II-A]{fedkls} is
\begin{align}
 s_{k,t,i}&=\Dkl{q_{t,0,i}}{q_{k,t,i}}\notag\\
 &=\frac{1}{2}\left[
 \frac{(\mu_{k,t,i}-\mu_{t,0,i})^2}{v_{k,t,i}}
 +\frac{v_{t,0,i}}{v_{k,t,i}}
 -1-\log\frac{v_{t,0,i}}{v_{k,t,i}}\right],
 \label{fkl:eq:score}
\end{align}
where $v$ denotes variance. Sort the scores so that $s_{k,t,(1)}\ge\cdots\ge s_{k,t,(d)}$. Let $\mathcal S_{k,t}(m)$ contain the indices of the largest $m$ scores, using a deterministic rule for ties. The server's reconstructed marginal is
\begin{equation}
 \widetilde q_{k,t,i}(m)=
 \begin{cases}
 q_{k,t,i},&i\in\mathcal S_{k,t}(m),\\
 q_{t,0,i},&i\notin\mathcal S_{k,t}(m).
 \end{cases}
 \label{fkl:eq:prior-fill}
\end{equation}

Define the discarded-score function
\begin{equation}
 D_{k,t}(m)=\sum_{j=m+1}^{d}s_{k,t,(j)},\qquad D_{k,t}(d)=0.
 \label{fkl:eq:discarded}
\end{equation}
It is nonincreasing since $D_{k,t}(m+1)=D_{k,t}(m)-s_{k,t,(m+1)}$. Factorization and~\eqref{fkl:eq:prior-fill} give
\begin{align}
 \Dkl{\widetilde q_{k,t}(m)}{q_{k,t}}
 &=\sum_{i=1}^{d}\Dkl{\widetilde q_{k,t,i}(m)}{q_{k,t,i}}\notag\\
 &=\sum_{i\notin\mathcal S_{k,t}(m)}s_{k,t,i}
 =D_{k,t}(m).
 \label{fkl:eq:exact-kl}
\end{align}
Thus the objective measures exact reconstruction KL relative to the computed local Gaussian distributions, assuming exact delivery of retained values. It is not an exact KL to the true Bayesian posterior, global loss, or test-accuracy degradation. Top-score retention minimizes this local reconstruction cost for a fixed count. Equivalently, minimizing discarded scores maximizes retained scores because the total score is fixed during allocation.

\subsubsection{Joint Problem and Exact Reduction}
The proposed problem is
\begin{equation}
 \begin{aligned}
 (\mathsf P_1)\quad
 \min_{\{m_{k,t},B_{k,t}\}}\quad&
 \sum_k\pi_k\Dkl{\widetilde q_{k,t}(m_{k,t})}{q_{k,t}}\\
 \mathrm{s.t.}\quad&
 L_{k,t}^{\UL}(m_{k,t})=U_tB_{k,t}c_{k,t},\quad\forall k,\\
 &\sum_kB_{k,t}\le B,\\
 &m_{k,t}\in\{1,\ldots,d\},\quad B_{k,t}>0.
 \end{aligned}
 \label{fkl:prob:joint}
\end{equation}
All $K$ clients participate. There is no whole-client exclusion variable and no allowed-ratio menu. Current local posteriors, aggregation weights, channel efficiencies, and preparation times are fixed inputs, not decision variables.

Define the required bandwidth in Hz by
\begin{equation}
 w_{k,t}(m)=\frac{L_{k,t}^{\UL}(m)}{U_tc_{k,t}}.
 \label{fkl:eq:bandwidth-cost}
\end{equation}
Substitution of~\eqref{fkl:eq:bandwidth-recovery} and~\eqref{fkl:eq:exact-kl} yields
\begin{equation}
 \begin{aligned}
 (\mathsf P_2)\quad \min_{\mathbf m_t}\quad&
 F_t(\mathbf m_t)=\sum_k\pi_kD_{k,t}(m_{k,t})\\
 \mathrm{s.t.}\quad&\sum_kw_{k,t}(m_{k,t})\le B,\\
 &m_{k,t}\in\{1,\ldots,d\},\quad\forall k.
 \end{aligned}
 \label{fkl:prob:counts}
\end{equation}

\begin{proposition}
For $U_t>0$, $B>0$, and $c_{k,t}>0$, problems $\mathsf P_1$ and $\mathsf P_2$ have identical feasible retained-count vectors and optimal reconstruction objectives. Every feasible $\widehat{\mathbf m}_t$ yields a feasible joint allocation through~\eqref{fkl:eq:bandwidth-recovery}.
\end{proposition}
\begin{proof}
The time equalities uniquely impose $B_{k,t}=w_{k,t}(m_{k,t})$. Their sum satisfies~\eqref{fkl:eq:bandwidth-budget} exactly when the reduced resource constraint holds. Conversely, assigning these bandwidths to any feasible count vector satisfies each time equality. Equation~\eqref{fkl:eq:exact-kl} gives the same objective in both problems.
\end{proof}
The same count-feasibility condition also characterizes a model with $T_{k,t}^{\UL}\le U_t$: any feasible bandwidth can be reduced to its required value when no additional lower bandwidth limits are imposed. Equality is the selected recovery policy, not a reason to restrict counts to special ratios.

\subsubsection{Feasibility and Endpoint Cases}
Require $U_t>0$ and positive channel efficiencies. Define
\begin{equation}
 W_t^{\min}=\sum_k\min\{w_{k,t}(1),w_{k,t}(d)\},\qquad
 W_t^{\dense}=\sum_kw_{k,t}(d).
 \label{fkl:eq:endpoint-costs}
\end{equation}
The minimum uses only these two endpoints because the sparse payload increases with $m$. If $W_t^{\min}>B$, the prescribed all-client problem is infeasible. If $W_t^{\dense}\le B$, full retention attains the global lower limit $F_t=0$. Otherwise, resource allocation is required, although zero scores may still permit a zero-distortion sparse solution.

The dense exception is material: a near-dense sparse message with $d+2bm\ge2bd$ is weakly dominated by the dense message, which uses no more bandwidth and has zero reconstruction cost. Consequently, the payload is not globally increasing over $m=1,\ldots,d$. All comparisons below use the exact piecewise expression~\eqref{fkl:eq:uplink-payload}.


\subsection{Lagrangian Decomposition and KL-Threshold Allocation}
\label{fkl:sec:solution}
We use only a shared-price Lagrangian method for the proposed allocator. Bandwidth elimination is exact, whereas the remaining integer allocation is approached through exact local price responses, a coordinated price search, and an explicitly feasible incumbent. General dual-decomposition and KKT concepts follow~\cite{decomposition,boyd}; the formulas below specialize them to the stated FedKLS objective.

\subsubsection{Separable Local Price Problems}
Suppress $t$ for this section. Introduce a multiplier $\lambda\ge0$ for the bandwidth constraint. The Lagrangian separates as
\begin{equation}
 \mathcal L(\mathbf m,\lambda)
 =\sum_k\left[\pi_kD_k(m_k)+\lambda w_k(m_k)\right]-\lambda B.
 \label{fkl:eq:lagrangian}
\end{equation}
For a broadcast price, each client returns
\begin{equation}
 m_k(\lambda)\in
 \argmin_{m\in\{1,\ldots,d\}}
 \underbrace{\left[\pi_kD_k(m)+\lambda w_k(m)\right]}_{F_k(m;\lambda)}.
 \label{fkl:eq:local-price}
\end{equation}
The client needs its own scores, $\pi_k$, $U$, and $c_k$, but not another client's posterior. The server coordinates the common price and checks the shared resource constraint. The price $\lambda$ is distinct from any regularization coefficient used in local Bayesian training.

\subsubsection{Exact Sparse Threshold and Dense Comparison}
For $1\le m\le d-2$, both $m$ and $m+1$ use the sparse format. The change in the local priced objective is
\begin{align}
 F_k(m+1;\lambda)-F_k(m;\lambda)
 &=-\pi_ks_{k,(m+1)}+\lambda\frac{2b}{Uc_k}.
 \label{fkl:eq:marginal-change}
\end{align}
Thus the next pair is worth retaining when
\begin{equation}
 s_{k,(m+1)}>\theta_k(\lambda),\qquad
 \theta_k(\lambda)=\frac{2b\lambda}{\pi_kUc_k}.
 \label{fkl:eq:threshold}
\end{equation}
Let $N_k(\lambda)=\#\{i:s_{k,i}>\theta_k(\lambda)\}$. Because scores decrease with rank, a sparse minimizer is
\begin{equation}
 m_k^{\spc}(\lambda)
 =\min\{d-1,\max\{1,N_k(\lambda)\}\}.
 \label{fkl:eq:sparse-response}
\end{equation}
Strict thresholding selects the lowest-count sparse minimizer when scores equal the threshold. A deterministic coordinate-index rule resolves score ties when constructing the actual message.

The full-retention option must be compared separately:
\begin{align}
 F_k^{\spc}(\lambda)
 &=\pi_kD_k(m_k^{\spc})
 +\lambda\frac{d+2b m_k^{\spc}}{Uc_k},\label{fkl:eq:sparse-price}\\
 F_k^{\dense}(\lambda)&=\lambda\frac{2bd}{Uc_k}.
 \label{fkl:eq:dense-price}
\end{align}
Choose the candidate with smaller priced objective; in a tie choose the smaller required bandwidth, then the smaller discarded KL if necessary. Equations~\eqref{fkl:eq:sparse-response}--\eqref{fkl:eq:dense-price} solve~\eqref{fkl:eq:local-price} exactly in mathematical arithmetic. In particular, the bitmap is included even though its constant cost cancels in~\eqref{fkl:eq:marginal-change}.

The resulting count is an integer computed from current scores, not an entry in a predefined ratio list. For fixed price, stronger channels and larger aggregation weights lower the threshold. Equivalently, the marginal benefit per additional Hz on the sparse branch is
\begin{equation}
 \rho_{k,i}=\frac{\pi_kUc_ks_{k,i}}{2b},
 \label{fkl:eq:priority}
\end{equation}
and the threshold retains marginal entries with $\rho_{k,i}>\lambda$.

\subsubsection{Server Price Search}
Define the total required bandwidth and the response objective by
\begin{equation}
 G(\lambda)=\sum_kw_k(m_k(\lambda)),\qquad
 F(\mathbf m(\lambda))=\sum_k\pi_kD_k(m_k(\lambda)).
 \label{fkl:eq:demand}
\end{equation}
For exact local minimizers, $G(\lambda)$ is nonincreasing, including dense-format switches. Indeed, let $m_1,m_2$ minimize one client's priced objective at $\lambda_1<\lambda_2$. Adding the two optimality inequalities gives
\begin{equation}
 (\lambda_2-\lambda_1)\big[w_k(m_2)-w_k(m_1)\big]\le0.
 \label{fkl:eq:monotonicity}
\end{equation}
Summing over clients proves the claim. The dense exception prevents replacing bandwidth demand by an unqualified count-only proxy; the server evaluates the actual piecewise payload at every query.

The server increases the price when $G(\lambda)>B$ and decreases it when a lower price may preserve more information while remaining feasible. It maintains a bracket with an infeasible low-price response and a feasible high-price response, and queries midpoint prices. An initial high price can be doubled until a feasible response is found, subject to a finite iteration and numerical-safety limit. At every query, the best feasible vector is retained rather than blindly accepting the last trial. A minimum-bandwidth vector from~\eqref{fkl:eq:endpoint-costs} provides a feasible fallback whenever the problem is feasible.

After choosing $\widehat{\mathbf m}$, the server grants
\begin{equation}
 \widehat B_k=\frac{L_k(\widehat m_k)}{Uc_k},\qquad
 \widehat r_k=\frac{\widehat m_k}{d}.
 \label{fkl:eq:final-allocation}
\end{equation}
Thus every idealized upload takes $U$ and the common completion target is met. Reallocating unused bandwidth to accelerate uploads would constitute a different completion policy, not the equality recovery in~\eqref{fkl:eq:final-allocation}.

\subsubsection{Dual Bound, Stopping Rule, and Optimality Scope}
The dual function is
\begin{equation}
 g(\lambda)=\sum_k\min_{1\le m\le d,\,m\in\mathbb Z}
 \left[\pi_kD_k(m)+\lambda w_k(m)\right]-\lambda B.
 \label{fkl:eq:dual}
\end{equation}
For any feasible incumbent $\widehat{\mathbf m}$ and $\lambda\ge0$,
\begin{equation}
 g(\lambda)\le F^\star\le F(\widehat{\mathbf m}),
 \label{fkl:eq:gap-bound}
\end{equation}
where $F^\star$ is the optimal value of~\eqref{fkl:prob:counts}. The first inequality follows because the Lagrangian is no greater than the primal objective at any feasible point, and $g$ minimizes that Lagrangian. The second follows from feasibility. No strong-duality assumption is needed.

Over the queried prices, record
\begin{equation}
 \mathrm{LB}=\max\{0,\max_jg(\lambda_j)\},\qquad
 \mathrm{UB}=F(\widehat{\mathbf m}).
 \label{fkl:eq:bounds}
\end{equation}
A requested reconstruction-objective accuracy is certified when
\begin{equation}
 \mathrm{UB}-\mathrm{LB}
 \le\varepsilon_{\mathrm{abs}}+
 \varepsilon_{\mathrm{rel}}\max\{1,\mathrm{UB}\}.
 \label{fkl:eq:gap-stop}
\end{equation}
A price-interval tolerance or iteration cap may also stop the algorithm, but does not by itself certify this objective accuracy. If exact local responses are feasible and satisfy
\begin{equation}
 \lambda\left(G(\lambda)-B\right)=0,
 \label{fkl:eq:complementarity}
\end{equation}
then $g(\lambda)=F(\mathbf m(\lambda))$, which certifies a global integer optimum for that instance.

In general, integer counts cause stepwise demands and can leave a nonzero duality gap. Bisection cannot recover every feasible integer combination. For illustration, after mandatory minimum allocations, suppose one optional pair costs two bandwidth units and reduces weighted KL by three, while another costs three units and reduces KL by four. A remaining budget of three units favors the second pair alone, but a common price selects the first pair whenever it selects the second. This illustrates why exact local thresholding does not automatically yield a globally optimal integer allocation. The separate dense format adds another nonconvex feature.

For a KKT interpretation, linearly interpolate $D_k$ between consecutive integer counts and restrict to the sparse branch. On $x\in[j,j+1]$,
\begin{equation}
 \overline D_k(x)=D_k(j)-(x-j)s_{k,(j+1)}.
 \label{fkl:eq:interpolation}
\end{equation}
Its slopes are nondecreasing, hence it is convex. Interior KKT stationarity for this continuous relaxation~\cite[Sec.~5.5]{boyd} gives
\begin{equation}
 0\in\pi_k\partial\overline D_k(x_k)+\lambda\frac{2b}{Uc_k}.
 \label{fkl:eq:kkt}
\end{equation}
At an interior integer breakpoint, this is
\begin{equation}
 \pi_ks_{k,(m_k)}\ge\frac{2b\lambda}{Uc_k}
 \ge\pi_ks_{k,(m_k+1)}.
 \label{fkl:eq:kkt-threshold}
\end{equation}
These conditions explain the threshold but do not make the original discrete problem convex. Accordingly, a returned feasible solution is denoted $\widehat{\mathbf m}$, not $\mathbf m^\star$ unless certified.

\begin{figure}[t]
\hrule\vspace{4pt}
\textbf{Procedure 1: Server-coordinated price allocation}
\vspace{3pt}
\footnotesize
\begin{algorithmic}[1]
\Require $U,B,c_k,\pi_k>0$; clients' sorted scores; finite query limits and tolerances.
\State Construct the minimum-bandwidth count vector.
\If{its required bandwidth exceeds $B$}
 \State \Return infeasible; do not force a positive count.
\EndIf
\State Initialize feasible incumbent $\widehat{\mathbf m}$; set $\mathrm{LB}=0$.
\State Also retain Version A's vector if feasible and better.
\If{$\sum_kw_k(d)\le B$}
 \State \Return all-dense counts and~\eqref{fkl:eq:final-allocation}.
\EndIf
\State Query exact client responses at $\lambda=0$.
\State Update bounds and the best feasible incumbent.
\If{the zero-price response is feasible}
 \State \Return its zero-distortion allocation.
\EndIf
\State Set $\lambda_{\mathrm{lo}}=0$, $\lambda_{\mathrm{hi}}>0$.
\State Query and double $\lambda_{\mathrm{hi}}$ until feasible or the bracketing limit is reached; record every query.
\If{a feasible high-price endpoint was found}
 \For{at most $J_{\max}$ bisection steps}
  \State Broadcast $\lambda=(\lambda_{\mathrm{lo}}+\lambda_{\mathrm{hi}})/2$.
  \State Clients solve~\eqref{fkl:eq:local-price} by thresholding and dense comparison; report count and discarded KL.
  \State Evaluate $G(\lambda),g(\lambda)$; update LB and incumbent.
  \If{$G(\lambda)>B$}
   \State $\lambda_{\mathrm{lo}}\gets\lambda$.
  \Else
   \State $\lambda_{\mathrm{hi}}\gets\lambda$.
  \EndIf
  \State Stop on~\eqref{fkl:eq:gap-stop}, price tolerance, or query limit.
 \EndFor
\EndIf
\State Optionally improve the incumbent by feasible additions or dense switches; recheck exact payloads.
\State \Return $\widehat m_k,\widehat r_k,\widehat B_k$, UB, LB, query count, and termination status.
\end{algorithmic}
\vspace{4pt}\hrule
\caption{Price coordination for unrestricted integer retention. The feasible incumbent is preserved even if the price-search limit is reached. No second optimization method is required by this procedure.}
\label{fkl:fig:algorithm}
\end{figure}

\subsubsection{Feasible Refinement and Computational Cost}
Procedure~1 may be followed by a bounded refinement that adds affordable posterior pairs according to~\eqref{fkl:eq:priority}, and checks switches to the dense format using their actual bandwidth difference. Accept only feasible changes that do not increase the reconstruction objective. Such additions can improve the incumbent, but are not a global-optimality proof and do not remove a duality gap in general. Previously valid dual bounds remain lower bounds after refinement.

A complete score sort and a reverse cumulative sum cost $O(d\log d)$ time and $O(d)$ storage at each client using a comparison sort. Afterwards, a binary search in sorted scores determines the threshold count in $O(\log d)$ time; tail evaluation and dense comparison are constant-time. Each price query uses $O(K)$ server arithmetic and a small response from each client, rather than uploading every coordinate score. With $J$ price queries, reporting scales as $O(KJ)$ scalar messages. This is a computational and signaling tradeoff, not a guarantee that repeated control exchanges have negligible elapsed time.


\subsection{Implementation and Evaluation Protocol}
\label{fkl:sec:implementation}
\subsubsection{Separation of Learning and Allocation}
Keep the existing local FOLA training and FedKLS reconstruction/aggregation intact. Introduce the wireless simulator and allocator as separate modules so that changing the resource policy does not silently change the local learner. A suggested interface is given in Table~\ref{fkl:tab:modules}. These are implementation contracts, not claims of completed training experiments.

\begin{table}[t]
\caption{Suggested implementation modules.}
\label{fkl:tab:modules}
\centering\footnotesize
\begin{tabular}{@{}>{\raggedright\arraybackslash}p{0.25\columnwidth}>{\raggedright\arraybackslash}p{0.68\columnwidth}@{}}
\toprule
Module & Responsibility and principal outputs\\
\midrule
Channel model & Distances, path gains, fading, PSD convention; current $c_{k,t}$ and multicast rate.\\
Client preparation & Full local training, positive posterior variances, KL scores, sorted indices, and tail sums.\\
Local price response & Given $\lambda,U_t,c_{k,t},\pi_k$, return exact priced count and $D_{k,t}(m)$.\\
Server allocator & Check feasibility; coordinate prices; maintain incumbent, bounds, and termination status.\\
Packet serializer & Send bitmap and top-$m$ pairs, or bitmap-free dense values; record actual byte count separately.\\
Round controller & Enforce the reporting barrier, grants, common upload interval, aggregation, and logging.\\
\bottomrule
\end{tabular}
\end{table}

For a simulated round, obtain the global posterior and downlink duration, execute each client's full local preparation, and compute $U_t$ using the maximum preparation duration. Refresh or generate channel conditions for the allocation/upload phase; do not assume an old channel measurement remains valid through arbitrarily long local training. The server then invokes Procedure~1. In an initial single-process simulator, client-response calls may be local function calls, but the server must not use the clients' full score vectors for decisions beyond the specified response interface. Account separately for the network exchanges those calls represent.

After a grant, client $k$ selects its stored top-$\widehat m_{k,t}$ indices and serializes the corresponding pairs. It does not train again during the price search. The server fills omitted marginals with the broadcast global values and preserves the original normalized weighted product~\cite[Sec.~II-A]{fedkls}:
\begin{equation}
 q_{t+1,0}(\boldsymbol\theta)
 \propto\prod_{k\in\mathcal K}
 \widetilde q_{k,t}(\boldsymbol\theta)^{\pi_k}.
 \label{fkl:eq:aggregation}
\end{equation}
Do not renormalize $\pi_k$ separately over only those clients that transmitted a particular coordinate; that would change the agreed prior-fill aggregation rule.

\subsubsection{Numerical and Feasibility Safeguards}
Use one source of truth for the piecewise payload in~\eqref{fkl:eq:uplink-payload}. Validate $U_t>0$, finite positive $c_{k,t}$, $1\le\widehat m_{k,t}\le d$, and
\begin{equation}
 \sum_k\widehat B_{k,t}\le B,\qquad
 \frac{L_k(\widehat m_{k,t})}{\widehat B_{k,t}c_{k,t}}=U_t
 \label{fkl:eq:postconditions}
\end{equation}
up to declared numerical tolerances. Never round an infeasible count upward to one or silently exclude a client. A violated all-client requirement must return an infeasibility status for the experiment controller to handle explicitly.

Compute score sums and objective comparisons in at least float64 for the initial numerical implementation. Construct tail sums by reverse accumulation rather than subtracting a nearly equal prefix from the total. Validate positive finite variances; apply the same damping/flooring convention to all compared learners. For variance ratio $v_0/v_k=1+u$, the variance contribution $u-\log(1+u)$ can be evaluated with a numerically stable logarithm near $u=0$. Reject nonfinite scores and distinguish tiny roundoff negatives from genuinely invalid inputs.

Use deterministic score ties and dense/sparse cost ties. A numerically approximate local response need not achieve the exact local Lagrangian minimum; its evaluated value is then not automatically a valid dual lower bound. Mathematical certificates in~\eqref{fkl:eq:gap-bound} require exact local minimization or rigorously bounded numerical errors. Without such error control, label the reported gap as a numerical diagnostic. Price iterations should have explicit limits and return the best feasible incumbent on termination.

The initial channel model fixes PSD, not total transmit power. Log $P_{k,t}=p_{k,t}^{\mathrm{PSD}}\widehat B_{k,t}$ and, under the equality timing policy, $E_{k,t}^{\UL}=P_{k,t}U_t$. If an experiment instead fixes total power, replace the rate and rederive the allocation cost; do not merely recompute SNR once and reuse the fixed-$c$ threshold.

\subsubsection{Equal-Bandwidth Baseline and Controlled Comparisons}
Use Version A as a baseline under the same reporting barrier and round target. Assign $B_{k,t}^{A}=B/K$ and define $Q_{k,t}=U_t(B/K)c_{k,t}$. Check $Q_{k,t}\ge2bd$ first, in which case $m_{k,t}^{A}=d$. Otherwise use
\begin{equation}
 m_{k,t}^{A}
 =\min\left\{d-1,
 \left\lfloor\frac{Q_{k,t}-d}{2b}\right\rfloor\right\},
 \label{fkl:eq:version-a}
\end{equation}
only when the result is at least one. Infeasibility under equal bandwidth does not imply infeasibility under joint allocation. The baseline transmits within $U_t$ and can finish early; the equality-based policy recovers its own bandwidths using~\eqref{fkl:eq:final-allocation}.

If Version A is feasible, its counts are feasible for the reduced joint problem because
\begin{equation}
 \sum_kw_{k,t}(m_{k,t}^{A})\le\sum_k B/K=B.
 \label{fkl:eq:baseline-incumbent}
\end{equation}
Initializing the feasible incumbent with this vector guarantees that the returned frozen-round reconstruction objective is no worse than Version A's, since it is never replaced by a worse feasible candidate. This statement compares identical local posteriors in one round; it does not guarantee higher accuracy or domination along separate training trajectories.

Start validation with frozen local posteriors. On small synthetic cases, compare the local threshold response against evaluation of all integer counts and compare the complete feasible solution against exhaustive enumeration of the tiny joint count space. Include zero scores, repeated scores, the dense exception, an all-dense feasible budget, an infeasible minimum budget, and heterogeneous $c_k$. Test that bandwidth demand is nonincreasing over increasing prices, that~\eqref{fkl:eq:postconditions} holds, and that valid lower and upper bounds bracket the tiny-instance optimum.

Then evaluate repeated federated training using matched initial models, local data partitions, channel draws, training settings, and round targets. Compare joint KL-aware allocation with feasible equal-bandwidth adaptation and a uniform-retention baseline. A random-selection ablation should match retained counts and assigned bandwidths to isolate the coordinate selector. Record accuracy and calibration versus cumulative time, reconstruction objective, count distributions, bandwidths, power/energy, price-query counts, feasibility failures, and objective gaps.

\subsubsection{Overhead Accounting and Scope of Claims}
The analytical round target omits price messages and server computation, but code should separately record or model both. Wall-clock time for in-memory function calls is not wireless-report latency. If their contribution is material, reserve a control-time allowance before allocation or explicitly subtract it from the upload budget. A finite price-query cap makes such an allowance operationally meaningful. Keep simulated parallel client durations separate from the sequential runtime of a single-machine simulation.

Under a fixed round duration, the intended benefit is better posterior preservation and potentially better learning within the same time budget, not an automatic reduction in per-round latency. The exact results are the reconstruction identity, bandwidth elimination, local priced solution, and conditional dual-bound certificate. The proposed price-based integer allocation is feasible by construction when the safeguards succeed, but is not universally globally optimal. Improved convergence, accuracy, calibration, and negligible control overhead remain hypotheses to be evaluated through the implementation.
